% Arquivo: latex/5-consideracoes.tex
% Seção 5: Considerações Finais

Esta revisão sistemática mapeou 34 estudos empíricos que avaliam os efeitos dos instrumentos da PNDR sobre indicadores socioeconômicos locais, cobrindo o período de 2005 a 2026. A análise do conjunto de evidências permite identificar cinco achados transversais.

Primeiro, há expressiva assimetria no volume de evidências entre instrumentos. Os Fundos Constitucionais concentram 28 dos 34 estudos, com predominância do FNE (24 estudos), enquanto os Fundos de Desenvolvimento (8 estudos) e os Incentivos Fiscais (9 estudos) permanecem substancialmente menos investigados. Essa desproporção é particularmente relevante dado o volume fiscal dos instrumentos menos avaliados: os Incentivos Fiscais da SUDAM e SUDENE totalizaram R\$ 126 bilhões em gastos tributários no período 2015--2024, sem que haja avaliações dedicadas aos incentivos da SUDAM ou análises de custo-efetividade de qualquer dos programas.

Segundo, os efeitos sobre o mercado de trabalho apresentam maior convergência entre os estudos do que os efeitos sobre o PIB \textit{per capita}. Para o emprego formal, há evidências consistentes de impacto positivo do FNE, com magnitude crescente nos primeiros anos pós-tratamento e efeitos mais pronunciados em micro e pequenas empresas e na modalidade de crédito para investimento. Para o PIB \textit{per capita}, os resultados são heterogêneos e fortemente condicionados à especificação econométrica, ao período amostral e à escala geográfica. Essa divergência sugere que os instrumentos podem ser mais eficazes na geração de postos de trabalho do que na elevação sustentada da renda média local.

Terceiro, a eficácia dos Fundos Constitucionais sobre o produto é condicionada ao nível inicial de renda dos municípios. Efeitos positivos concentram-se em faixas intermediárias de desenvolvimento, são nulos ou inconclusivos em municípios de baixa renda e decrescentes em municípios de alta renda. Esse padrão é consistente com a hipótese de capacidade de absorção: municípios mais pobres podem carecer do capital humano e da infraestrutura necessários para converter crédito subsidiado em crescimento sustentado, enquanto municípios de renda intermediária possuem o mínimo de condições para ativar o efeito multiplicador dos investimentos.

Quarto, os Fundos de Desenvolvimento apresentam indícios de impacto expressivo sobre o PIB \textit{per capita} --- entre 19\% e 24\% ---, com presença de efeitos de transbordamento espacial para municípios vizinhos. Esses valores são consideravelmente superiores aos estimados para os Fundos Constitucionais, o que pode refletir a escala dos empreendimentos financiados e sua concentração em infraestrutura e indústria de transformação. A base empírica, contudo, restringe-se a estudos recentes (2024--2026) com metodologias semelhantes e períodos parcialmente sobrepostos.

Quinto, a variável salário médio não apresenta variação significativa na maioria das avaliações dos instrumentos federais, indicando que novos postos de trabalho são criados ao nível salarial vigente, sem ganhos adicionais de produtividade. No âmbito estadual, a avaliação do Prodepe (PE) aponta redução salarial concomitante ao aumento do emprego, suscitando preocupação sobre a qualidade dos postos gerados por incentivos subnacionais.

Esses achados têm implicações para o desenho da política regional. A constatação de que crédito subsidiado é insuficiente para dinamizar municípios de baixa renda sugere a necessidade de complementaridade com investimentos em infraestrutura, capital humano e capacidades institucionais locais --- condições prévias à eficácia do instrumento. A ausência de ganhos salariais associados à geração de empregos indica que os instrumentos, em sua configuração atual, tendem a expandir a atividade econômica sem transformar a estrutura produtiva. A concentração dos Incentivos Fiscais em municípios de maior dinamismo, documentada na literatura, contrasta com o objetivo redistributivo da PNDR e aponta para a necessidade de revisão dos critérios de concessão.

Do ponto de vista metodológico, este artigo constitui a primeira revisão sistemática dedicada aos instrumentos de financiamento da PNDR. A estratégia de busca abrangeu cinco bases de dados acadêmicas, a classificação dos registros empregou modelo de linguagem de grande escala com revisão humana integral, a seleção dos estudos foi respaldada por índice de citação cruzada --- que permitiu justificar a inclusão de literatura cinzenta com base em reconhecimento pela própria comunidade acadêmica --- e a deduplicação percorreu quatro fases complementares (DOI exato, similaridade de título, identidade de PDF e verificação manual de textos para discussão e versões de congresso).

Não obstante, o estudo possui limitações que devem ser explicitadas. A heterogeneidade de métricas, períodos amostrais e especificações econométricas entre os estudos primários inviabiliza a condução de meta-análise formal, restringindo a síntese ao plano qualitativo. Os resultados desta revisão dependem, ademais, da qualidade e da comparabilidade dos trabalhos incluídos, cujas estratégias de identificação causal variam substancialmente em robustez. Embora a inclusão de literatura cinzenta atenue o risco de viés de publicação, tal risco não pode ser inteiramente eliminado. A triagem assistida por modelo de linguagem demandou correção em dois terços dos registros classificados, o que evidencia a necessidade de validação integral pelo pesquisador e impede a delegação autônoma dessa etapa. Por fim, a busca restringiu-se a cinco bases de dados, de modo que estudos indexados exclusivamente em outros repositórios podem não ter sido capturados.

Do ponto de vista da agenda de pesquisa, três prioridades emergem desta revisão: (i)~a ampliação das evidências sobre Fundos de Desenvolvimento --- especialmente FDA e FDCO, para os quais inexistem avaliações quantitativas dedicadas --- e sobre os Incentivos Fiscais da SUDAM; (ii)~a condução de estudos que avaliem simultaneamente os três instrumentos e suas interações, dado que a quase totalidade das avaliações examina cada instrumento isoladamente, sem capturar possíveis sinergias ou redundâncias; e (iii)~a adoção de desenhos quase experimentais com maior poder de identificação causal --- como regressão com descontinuidade e controle sintético ---, que permanecem sub-representados na literatura frente a métodos com maior fragilidade na separação do efeito da política de fatores confundidores.
